问题四在问题二、三的已验收框架上引入附件 4 的交付时段时变电价。难点不再只是负荷与光伏偏差，同一物理缺口在不同时段还有不同价格权重：把价格当作 0:00 已知常数会错估“提前购电—储能转移—后续释放”的价值；假定 0:00 已知全天实际价同样不可执行。

因此，问题四的核心变化是因果价格信息，而不是重新放宽问题二、问题三已经建立的负荷与光伏非前瞻性边界。波动电价下的日前额度策略（Q4-2）在联合残差日前评估与逐十分钟 MPC 上，把历史价格残差与负荷、光伏残差同步配对；波动电价下的多时点调整策略（Q4-3）沿用全时点价值触发规则，把调整结算与紧急购电改为按交付时段实际价格计价。两条链共享同一价格预测模块，但分别对应题设的两类结果。

与问题二、问题三相同，1 月仍只作预测历史与 SOC 预热，题设输出区间为 2--12 月。两类策略均在相同设备边界、信息时序和输出区间下核算，避免把预热成本或不同价格机制下的差额误作策略优劣。
